%&bericht

%%%%%%%%%%%%%%%%%%%%%%%%%%%%%%%%%%%%%%%%%%%%%%%%%%%%%%%%%%%%%%%%%%%%%%%%%%%%%%%
%% Descr:       Vorlage für Berichte der DHBW-Karlsruhe
%% Author:      Prof. Dr. Jürgen Vollmer, juergen.vollmer@dhbw-karlsruhe.de
%% $Id: bericht.tex,v 1.25 2020/03/13 15:07:45 vollmer Exp $
%%  -*- coding: utf-8 -*-
%%%%%%%%%%%%%%%%%%%%%%%%%%%%%%%%%%%%%%%%%%%%%%%%%%%%%%%%%%%%%%%%%%%%%%%%%%%%%%%

\documentclass[
   ngerman          % neue deutsche Rechtschreibung
  ,a4paper 
  ,oneside         % Papiergrösse
% ,twoside          % Zweiseitiger Druck (rechts/links)
% ,10pt             % Schriftgrösse
% ,11pt
  ,12pt
  ,pdftex
  ,parskip=half
  ,headsepline
  ,abstraction
  ,footsepline  
%  ,disable         % Todo-Markierungen auschalten
]{report}
\usepackage[ngerman]{babel}
\usepackage[T1]{fontenc}
% Bitte die Codierung Ihrer Dateien auswählen:
% \usepackage[latin1]{inputenc}    % Für UNIX mit ISO-LATIN-codierten Dateien
% \usepackage[applemac]{inputenc}  % Für Apple Mac
% \usepackage[ansinew]{inputenc}   % Für Microsoft Windows
\usepackage[utf8]{inputenc}        % UTF-8 codierte Dateien
                                   % Dieses Dokument ist unter Unix erstellt, daher
                                   % wird diese Input-Codierung benutzt.

\usepackage{bericht}
\usepackage{listings}


%% ACHTUNG, wenn man eine eigene Formatdatei (bericht.fmt) benutzt, werden Änderungen an bericht.sty
%% erst wirksam, wenn die Format-Datei neu erzeugt wurde!!!
%% Genauer alle Änderungen, die textuell vor der nächsten Zeile ".... endofdump...." stehen
%% werden erst wirksam, wenn die Formatdatei neu erzeugt wurde
\csname endofdump\endcsname

%%%%%%%%%%%%%%%%%%%%%%%%%%%%%%%%%%%%%%%%%%%%%%%%%%%%%%%%%%%%%%%%%%%%%%%%%%%%%%%
%% Angaben zur Arbeit
%%%%%%%%%%%%%%%%%%%%%%%%%%%%%%%%%%%%%%%%%%%%%%%%%%%%%%%%%%%%%%%%%%%%%%%%%%%%%%%

\newcommand{\Autor}{Matthias Schuhmacher und Marlene Rieder}
\newcommand{\MatrikelNummer}{4128647 und 8261867}
\newcommand{\Kursbezeichnung}{tinf21b3 und tinf21b5}

%\newcommand{\FirmenName}{Siemens AG}
%\newcommand{\FirmenStadt}{Karlsruhe}
%\newcommand{\FirmenLogoDeckblatt}{{\includegraphics[width=5cm]{Siemens_Logo_transp}}}

% Falls es kein Firmenlogo gibt:
%  \newcommand{\FirmenLogoDeckblatt}{}

%\newcommand{\BetreuerFirma}{Boris Rojs}
\newcommand{\BetreuerDHBW}{Prof. Dr. Marcus Strand}

%%%%%%%%%%%%%%%%%%%%%%%%%%%%%%%%%%%%%%%%%%%%%%%%%%%%%%%%%%%%%%%%%%%%%%%%%%%%%%%%%%%%%

% Wird auf dem Deckblatt und in der Erklärung benutzt:
%\newcommand{\Was}{Projekt-/Studien-/Bachleorarbeit}
%\newcommand{\Was}{Projektarbeit}
\newcommand{\Was}{Studienarbeit}
%\newcommand{\Was}{Bachleorarbeit}

%%%%%%%%%%%%%%%%%%%%%%%%%%%%%%%%%%%%%%%%%%%%%%%%%%%%%%%%%%%%%%%%%%%%%%%%%%%%%%%%%%%%%

\newcommand{\Titel}{Pepper VR – Teleoperation eines humanoiden Roboter auf Basis der Analyse menschlicher Bewegung}
\newcommand{\AbgabeDatum}{20. Mai 2024}

\newcommand{\Dauer}{300 Stunden}

% \newcommand{\Abschluss}{Bachelor of Engineering}
\newcommand{\Abschluss}{Bachelor of Science}

%\newcommand{\Studiengang}{Informatik / Informationstechnik}
 \newcommand{\Studiengang}{Informatik / Angewandte Informatik}

\hypersetup{%%
  pdfauthor={\Autor},
  pdftitle={\Titel},
  pdfsubject={\Was}
}

%%%%%%%%%%%%%%%%%%%%%%%%%%%%%%%%%%%%%%%%%%%%%%%%%%%%%%%%%%%%%%%%%%%%%%%%%%%%%%%

% Wenn \includeonly{..} benutzt wird, werden nur diese Kaptitel ausgegeben.
%\includeonly{
 % abk
%  Einleitung
 %,Grundlagen
 %,Die Projekte
 %,Reflexion und Ausblick
%,changelog%
%}

%%%%%%%%%%%%%%%%%%%%%%%%%%%%%%%%%%%%%%%%%%%%%%%%%%%%%%%%%%%%%%%%%%%%%%%%%%%%%%%

% Benutzt man das "biblatex"-Paket, dann muß das hier stehen:
% siehe auch die mit BIBLATEX markierten Zeilen in bericht.sty
\bibliography{bericht.bib}
\def\SymbReg{\textsuperscript{\textregistered}}
\nocite{*}
\begin{document}

%%%%%%%%%%%%%%%%%%%%%%%%%%%%%%%%%%%%%%%%%%%%%%%%%%%%%%%%%%%%%%%%%%%%%%%%%%%%%%%

\begin{titlepage}
\begin{center}
\vspace*{-2cm}
\includegraphics[width=4cm]{Abbildungen/dhbw-logo.png}\\[2cm]
{\Huge \Titel}\\[1cm]
{\Huge\scshape \Was}\\[1cm]
{\large für die Prüfung zum}\\[0.5cm]
{\Large \Abschluss}\\[0.5cm]
{\large des Studienganges \Studiengang}\\[0.5cm]
{\large an der}\\[0.5cm]
{\large Dualen Hochschule Baden-Württemberg Karlsruhe}\\[0.5cm]
{\large von}\\[0.5cm]
{\large\bfseries \Autor}\\[1cm]
{\large Abgabedatum \AbgabeDatum}
\vfill
\end{center}
\begin{tabular}{l@{\hspace{2cm}}l}
Bearbeitungszeitraum	         & \Dauer 			\\
Matrikelnummer	                 & \MatrikelNummer		\\
Kurs			         & \Kursbezeichnung		\\
%Ausbildungsfirma	         & \FirmenName			\\
%			         & \FirmenStadt			\\
%Betreuer der Ausbildungsfirma	 & \BetreuerFirma		\\
Gutachter der Studienakademie	 & \BetreuerDHBW		\\
\end{tabular}
\end{titlepage}

%%%%%%%%%%%%%%%NEWCOMMANDS_EXT%%%%%%%%%%%%%%%%%%%%
%Simple way to include images beside the Text
%  #1 right/left (R/L)
%  #2 Caption / filename
%  #3 Size in cm
%  #4 Vertical offset (negative possible)
%  #5 Source
\def\wrapfig #1#2#3#4#5{
    \StrSubstitute{#2}{-}{\medspace}[\text]
    \begin{wrapfigure}{#1}{#3cm}
        \vspace{#4cm}
        \includegraphics[width=#3cm]{Abbildungen/#2}
        \vspace{-1cm}
        \caption{\text \if#5- \else \protect\cite{#5} \fi }
        \label{#2}
    \end{wrapfigure}
}

%Simple inclue for default images
%  #1 Caption / filename
%  #2 Size in cm
%  #3 Source
\def\normfig #1#2#3{
    \StrSubstitute{#1}{-}{\medspace}[\text]
    \begin{figure}[H]
        \centering
        \includegraphics[width=#2\textwidth]{Abbildungen/#1}
        \caption{\text \if#3- \else \protect\cite{#3} \fi }
        \label{#1}
    \end{figure}
}

%Simple Listing
%  #1 Caption
%  #2 label / filename
%  #3 Language (has to fit style)
\def\code #1#2{
    \begin{figure}[!ht]
        \begin{framed}
            \lstinputlisting[language=c++, caption=#1, label=#2]{code/#2}
        \end{framed}
    \end{figure}
}



%%%%%%%%%%%%%%%%%%%%%%%%%%%%%%%%%%%%%%%%%%%%%%%%%%%%%%%%%%%%%%%%%%%%%%%%%%%%%%%

%%%%%%%%%%%%%%%%%%%%%%%%%%%%%%%%%%%%%%%%%%%%%%%%%%%%%%%%%%%%%%%%%%%%%%%%%%%%%%%

%% Descr:       Vorlage für Berichte der DHBW-Karlsruhe, Erklärung

%% Author:      Prof. Dr. Jürgen Vollmer, vollmer@dhbw-karlsruhe.de

%% $Id: erklaerung.tex,v 1.11 2020/03/13 14:24:42 vollmer Exp $

%% -*- coding: utf-8 -*-

%%%%%%%%%%%%%%%%%%%%%%%%%%%%%%%%%%%%%%%%%%%%%%%%%%%%%%%%%%%%%%%%%%%%%%%%%%%%%%%

 

% In Bachelorarbeiten muss eine schriftliche Erklärung abgegeben werden.

% Hierin bestätigen die Studierenden, dass die Bachelorarbeit, etc.

% selbständig verfasst und sämtliche Quellen und Hilfsmittel angegeben sind. Diese Erklärung

% bildet das zweite Blatt der Arbeit. Der Text dieser Erklärung muss auf einer separaten Seite

% wie unten angegeben lauten.

 

\newpage
\thispagestyle{empty}
\begin{framed}
\begin{center}
\Large\bfseries Erklärung
\end{center}
\medskip
\noindent
% siehe §5(3) der \enquote{Studien- und Prüfungsordnung DHBW Technik} vom 29.\,9.\,2017 und Anhang 1.1.13
Ich versichere hiermit, dass ich meine \Was mit dem Thema:
\enquote{\Titel}
selbstständig verfasst und keine anderen als die angegebenen Quellen und Hilfsmittel benutzt habe. Ich versichere zudem, dass die eingereichte elektronische Fassung mit der gedruckten Fassung übereinstimmt.
\vspace{3cm}
\\
\noindent
\underline{\hspace{4cm}}\hfill\underline{\hspace{6cm}}\\
Ort~~~~~Datum\hfill Unterschrift\hspace{4cm}
\end{framed} 

\vfill
\begin{framed}
\begin{center}
\Large\bfseries Sperrvermerk
\end{center}
\medskip
\noindent
Die vorliegende Projektarbeit T3\_2000 mit dem Titel
\enquote{\Titel}
enthält unternehmensinterne bzw. vertrauliche Informationen der Atruvia AG, ist deshalb mit einem Sperrvermerk versehen und wird ausschließlich zu Prüfungszwecken am Studiengang \Studiengang der Dualen Hochschule Baden-Württemberg Karlsruhe vorgelegt.
\newline
Sie ist ausschließlich zur Einsicht durch den zugeteilten Gutachter, die Leitung des Studiengangs und ggf. den Prüfungsausschuss des Studiengangs bestimmt. Es ist untersagt,
\newline
\indent - den Inhalt dieser Arbeit (einschließlich Daten, Abbildungen, Tabellen, Zeichnungen
\indent   usw.) als Ganzes oder auszugsweise weiterzugeben. Kopien oder Abschriften dieser
\indent   Arbeit (einschließlich Daten, Abbildungen, Tabellen, Zeichnungen usw.) als Ganzes
\indent   oder auszugsweise anzufertigen.
    \newline    
\indent - diese Arbeit zu veröffentlichen bzw. digital, elektronisch oder virtuell zur Verfügung zu
\indent   stellen.
    \newline
Jede anderweitige Einsichtnahme und Veröffentlichung - auch von Teilen der Arbeit - bedarf der vorherigen Zustimmung durch den Verfasser und die Atruvia AG.

\end{framed}

%%%%%%%%%%%%%%%%%%%%%%%%%%%%%%%%%%%%%%%%%%%%%%%%%%%%%%%%%%%%%%%%%%%%%%%%%%%%%%%
\endinput
%%%%%%%%%%%%%%%%%%%%%%%%%%%%%%%%%%%%%%%%%%%%%%%%%%%%%%%%%%%%%%%%%%%%%%%%%%%%%%%


%%%%%%%%%%%%%%%%%%%%%%%%%%%%%%%%%%%%%%%%%%%%%%%%%%%%%%%%%%%%%%%%%%%%%%%%%%%%%%%

\begin{abstract}
Das Ziel der vorliegenden Studienarbeit ist es, eine mögliche Verbindung zwischen einer Virtual-Reality Brille und dem humanoiden Roboter Pepper herzustellen. Das Kamerabild des Roboters soll auf der Brille angezeigt werden, ebenfalls soll es möglich sein, den Roboter mit Hilfe der Controller der Brille zu steuern.
\end{abstract}

\newpage
\tableofcontents           % Inhaltsverzeichnis hier ausgeben
\listoffigures             % Liste der Abbildungen
\listoftables              % Liste der Tabellen
\lstlistoflistings         % Liste der Listings
\listofequations           % Liste der Formeln

% Jetzt kommt der "eigentliche" Text
%%%%%%%%%%%%%%%%%%%%%%%%%%%%%%%%%%%%%%%%%%%%%%%%%%%%%%%%%%%%%%%%%%%%%%%%%%%%%%
%% Descr:       Vorlage für Berichte der DHBW-Karlsruhe, Datei mit Abkürzungen
%% Author:      Prof. Dr. Jürgen Vollmer, vollmer@dhbw-karlsruhe.de
%% $Id: abk.tex,v 1.4 2017/10/06 14:02:03 vollmer Exp $
%% -*- coding: utf-8 -*-
%%%%%%%%%%%%%%%%%%%%%%%%%%%%%%%%%%%%%%%%%%%%%%%%%%%%%%%%%%%%%%%%%%%%%%%%%%%%%%%

\chapter*{Abkürzungsverzeichnis}                   % chapter*{..} -->   keine Nummer, kein "Kapitel"
						         % Nicht ins Inhaltsverzeichnis
% \addcontentsline{toc}{chapter}{Akürzungsverzeichnis}   % Damit das doch ins Inhaltsverzeichnis kommt

% Hier werden die Abkürzungen definiert
\begin{acronym}[DHBW]
  % \acro{Name}{Darstellung der Abkürzung}{Langform der Abkürzung}
  \acro{Abk}[Abk.]{Abkürzung}
  \acro{DHBW} {Duale Hochschule Baden-Württemberg}
  \acro{SAS}[SAS]{Société par actions simplifiée}
  \acro{B2B}[B2b]{Business-to-Business}
  \acro{TH}[TH]{Technische Hochschule}
  \acro{API}[API]{Application Programming Interface}
  \acro{SDK}[SDK]{Software Development Kit}
  \acro{URL}[URL]{Uniform Resource Locator}
  \acro{IP}[IP]{Internet Protocol}
  \acro{RPC}[RPC]{Remote Procedure Call}
  \acro{LPC}[LPC]{Local Procedure Call}
  
  % Folgendes benutzen, wenn der Plural einer Abk. benöigt wird
  % \newacroplural{Name}{Darstellung der Abkürzung}{Langform der Abkürzung}
  \newacroplural{Abk}[Abk-en]{Abkürzungen}
  % Wenn neicht benutzt, erscheint diese Abk. nicht in der Liste
  \acro{NUA}{Not Used Acronym}
\end{acronym}
              % Abkürzungsverzeichnis

\chapter{Einleitung}
Im Umfeld der \ac{DHBW} Karlsruhe werden immer wieder Projekte im Bereich der Robotik an Studierende im Rahmen verschiedener Arbeiten vergeben. Ebenso dieses Projekt, welches im Rahmen der Studienarbeit an die beiden Studierenden Matthias Schuhmacher und Marlene Rieder übergeben wurde. 
\\

\noindent
Im Umfeld der \ac{DHBW} Karlsruhe werden immer wieder Projekte im Bereich der Robotik an Studierende im Rahmen verschiedener Arbeiten vergeben. Ebenso dieses Projekt, welches im Rahmen der Studienarbeit an die beiden Studierenden Matthias Schuhmacher und Marlene Rieder übergeben wurde.
\\

\noindent 
In unserer immer digitaler werdenden Welt spiele sowohl Robotik als auch \ac{VR} eine immer größer werdende Rolle, diese zwei Welten sollen durch das, im folgenden beschriebene Projekt verbunden werden.
\\

\noindent
Ziel dieses Projektes ist es den humanoiden Roboter Pepper mit Hilfe der \ac{VR} Brille Meta Quest 3 zu steuern. Dies soll ermöglicht werden, indem das Kamerabild des Roboters an die Brille übertragen wird und sich die Bewegungen des Roboters durch die Controller der Brille steuern lassen. Der Nutzer soll also sehen, was sich vor dem Roboter befindet und ihn dann steuern können.

\section{Motivation}
Die Teleoperation humanoider Roboter bietet eine Vielzahl von Vorteilen. Durch die Fernsteuerung können Roboter in gefährlichen oder unzugänglichen Umgebungen eingesetzt werden, was die Sicherheit der menschlichen Bediener erhöht. Außerdem ermöglicht die Verwendung von VR-Technologie eine intuitivere und natürlichere Steuerung der Roboter, da die Bediener das Gefühl haben, direkt vor Ort zu sein und in Echtzeit mit der Umgebung zu interagieren.

\section{Ziele und Fragestellungen}
Das Hauptziel dieser Arbeit ist es, eine funktionsfähige Verbindung zwischen der MetaQuest 3 VR-Brille und dem humanoiden Roboter Pepper herzustellen. Dabei sollen folgende Fragestellungen untersucht werden:
\begin{itemize}
    \item Wie kann die Datenübertragung zwischen der VR-Brille und dem Roboter optimiert werden, um eine latenzfreie Steuerung zu ermöglichen?
    \item Welche Herausforderungen ergeben sich bei der Integration von Unity und ROS zur Steuerung des Roboters?
    \item Welche technischen und ergonomischen Anforderungen müssen bei der Entwicklung der VR-Steuerung berücksichtigt werden?
\end{itemize}

\section{Methodik}
Zur Beantwortung dieser Fragestellungen wird ein methodischer Ansatz verfolgt, der sowohl theoretische als auch praktische Aspekte umfasst. Zunächst wird eine umfassende Literaturrecherche durchgeführt, um den aktuellen Stand der Technik und die verfügbaren Technologien zu analysieren. Anschließend wird ein Prototyp entwickelt, der die Verbindung zwischen der VR-Brille und dem Roboter herstellt. Dieser Prototyp wird in mehreren Iterationen getestet und optimiert, um die bestmögliche Leistung zu erzielen.

\section{Struktur der Arbeit}
Diese Arbeit gliedert sich wie folgt:
\begin{itemize}
    \item Kapitel 2 beschreibt die theoretischen Grundlagen und den aktuellen Stand der Technik.
    \item Kapitel 3 erläutert die technischen Spezifikationen der verwendeten Hardware und Software.
    \item Kapitel 4 stellt den Entwicklungsprozess und die Implementierung des Prototyps dar.
    \item Kapitel 5 präsentiert die Testergebnisse und Diskussion.
    \item Kapitel 6 gibt eine Zusammenfassung und Ausblick auf zukünftige Arbeiten.
\end{itemize}


\chapter{Grundlagen}\label{chap:Grundlagen}

\section{Pepper}\label{sec:Pepper}
Pepper ist ein humanoider Roboter, der entwickelt wurde, um die Gefühle und Gesten von Menschen zu analysieren und basierend auf diesen, darauf zu reagieren. Das Projekt entstand durch eine Zusammenarbeit des französischen Unternehmens Aldebaran Robotics \ac{SAS} und des japanischen Telekommunikations- und Medienkonzerns SoftBank Mobile Corp. Ziel diese Projektes war es, einen humanuiden ``Roboter-Gefährten'' oder einen ``persönlichen Roboter-Freund'' zu schaffen, der zunächst im Gewerbesektor in Verkaufsräumen, an Empfangstischen oder in Bildungs- und Gesundheitseinrichtungen eingesetzt werden sollte. Die Produktion wurde jedoch aufgrund geringer Nachfrage bis auf Weiteres pausiert.\\
Das Konzept von Pepper distanziert sich von herkömmlichen Industrierobotern und reinen Spielzeugrobotern, indem er als informativer und kommunikativer Begleiter konzipiert wurde. Sein Aussehen, das im etwa an die Größe eines Kindes angelehnt ist, sowie ein freundliches Gesicht und eine kindliche Stimme sind im ästhetischen Konzept von ``kawaii'' (japanisch für ``niedlich'' oder auch ``liebenswert'') gehalten.\\
Pepper wurde im Rahmen einer Präsentation am 5. Juni 2014 als der ``erste persönliche Roboter der Welt mit Emotionen'' vorgestellt. Die Vermarktung begann damit, dass SoftBank Pepper-Geräte in ihren Verkaufsräumen einsetzte, um Kunden zu unterhalten und zu informieren. Die Roboter sollte dabei den Umgang mit Kunden erlernen, um zukünftige Anwendungsmöglichkeiten zu erforschen. Verkauft wurde offiziell ab dem 3. Juli 2015 zu einem Preis von 198.000 Yen pro Einheit, zuzüglich monatlicher Gebühren für Zusatzleistungen. Im Laufe der Zeit wurde Pepper auch für den Einsatz in weiteren Unternehmen und Einrichtungen verfügbar gemacht.\\
Pepper wird mit einer Grundausstattung an Anwendungen geliefert, jedoch sind für spezifische Anwendungen, individuell entwickelte Softwarelösungen erforderlich wie auch zum Beispiel in diesem. SoftBank ermöglichte unabhängigen Entwicklern durch die Veröffentlichung der Schnittstellen den Zugang zu einem Interface für Applikationsprogramme, um zusätzliche Anwendungen für Pepper zu erstellen. Das NAOqi-Framework welches für diesen Nutzen bereitgestellt wurde, beinhaltet eine \ac{API}, eine \ac{SDK} und weitere Tools, welche in den Sprachen Python und C++ uneingeschränkten Zugriff auf die Komponenten, Sensoren und Aktoren des Roboters bieten, dazu später Ausführlicheres in \autoref{sec:NAOqi}. Mit Hilfe diese Interfaces haben verschiedene Unternehmen integrierte Lösungen entwickelt, die Pepper beispielsweise bei der Kundenberatung unterstützen können.\\
Das Design von Pepper ist dem Menschen ähnlich und umfasst einen Kopf mit integrierten Mikrofonen und Kameras sowie einen Torso mit weiteren Sensoren für Stabilität und Sicherheit. Der Roboter verfügt über verschiedene Mechaniken, die es ihm ermöglichen, sich flüssig zu bewegen und mit Personen zu interagieren. Durch die Verwendung von Kameras und bereitgestellter Software ist Pepper in der Lage, Emotionen bei seinen Gesprächspartnern zu erkennen und darauf zu reagieren, obwohl er selbst keine Mimik besitzt. Sicherheitsvorkehrungen wie Abstandssensoren und Stabilisatoren gewährleisten einen sicheren Einsatz von Pepper in verschiedenen Umgebungen. Diese können jedoch bedingt durch den Entwickler deaktiviert werden, um den Roboter in komplexeren oder laborähnlichen Umgebungen zu betreiben.\\

\section{Publish and Subscribe}\label{sec:PublishAndSubscribe}
Bevor wir uns mit den Technologien und Herangehensweisen des Pepper Roboters beziehungsweise dessen Betriebssystem beschäftigen, müssen wir uns mit dem Publish-Subscribe-Modell auseinandersetzen.\\
Das Publish-Subscribe-Modell ist ein Paradigma für einen effektiven Nachrichtenaustausch in verteilten Systemen. Erstmals publiziert in einem Paper der \ac{ACM} von 1987\cite{wiki_publish_subscribe_pattern}, ermöglicht es die flexible Kommunikation zwischen verschiedenen Komponenten, indem es einen Mechanismus bereitstellt, über den Nachrichten von einem Sender, dem Publisher, an einen oder mehrere Empfänger, den Subscribern, verteilt werden können. Ein zentrales Element dieses Modells ist dabei Nachrichtenbroker oftmals auch nur als Borker referenziert, der als Vermittler zwischen Publishern und Subscribenden fungiert.\\
Der Broker empfängt Nachrichten vom Publisher und organisiert sie in verschiedene Kategorien, welche als Topics bezeichnet werden. Diese Topics dienen als thematische Selektionsmerkmale, die den Inhalt der Nachrichten beschreiben. Sie können in den verschiedensten Datentypen erfolgen, einfache Ganzzahlen, Texte bis hin zu Bildern oder weitaus komplexeren Datenstrukturen. Durch diese Kategorisierung in den Topics, wird eine gezielte Auswahl und Weiterleitung der Nachrichten ermöglicht.\\
Subscriber können bestimmte Topics abonnieren, die ihren Gewünschten entsprechen. Dadurch erhalten sie nur die Nachrichten, die zu den von ihnen gewählten Themen gehören. Dieser Prozess der Nachrichtenauswahl und -verarbeitung wird als Filterung bezeichnet und kann auf zwei Arten erfolgen: themenbasiert und inhaltsbasiert.\\
Im themenbasierten Ansatz erhalten Abonnenten alle Nachrichten zu den Topics, welche sie abonniert haben. Der Publisher definiert die verfügbaren Topics, aus welchen die Abonnenten selbst wählen können. Im inhaltsbasierten Ansatz dagegen, werden Nachrichten nur an die Abonnenten weitergeleitet, wenn sie den vom dessen festgelegten Kriterien entsprechen. Dabei ist der Abonnent für die Spezifikation dieser Kriterien verantwortlich\cite{itwissen_publish_subscribe_model}.\\
Das Publish-Subscribe-Modell bietetdurch seine Herangehensweise eine hohe Flexibilität und Skalierbarkeit, und wird daher in verschiedenen Anwendungsbereichen eingesetzt. Es ermöglicht eine Entkopplung von Nachrichtenerzeugung und -verarbeitung, was gerade in der Entwicklung verteilter Systeme Prozesse erleichtert. Die vermutlich bekannteste Implementierung des Modell ist das MQTT-Protokoll, welches konzipiert wurde um Telemetriedeaten zwischen Sensoren und Servern zu übertragen speziell in unzuverlässigen Umgebungen\cite{elektronik_kompendium_publish_subscribe}.\\
Weiter findet das Publish-Subscribe-Modell Anwendung in einer Vielzahl von weiteren Systemen und Technologien. Beispiele dafür sind der \ac{JMS}, der \ac{DDS} oder wie im Fall diese Projektes die \ac{MOM} zu denen auch \ac{ROS} gehört.\\

\section{NAOqi}\label{sec:NAOqi}
Im folgenden Abschnitt wird das NAOqi-Framework, welches auf dem Pepper Roboter läuft, genauer erläutert.\\

\subsection{Was ist NAOqi?}\label{subsec:WasIstNAOqi}
NAOqi ist die Bezeichnung für die Hauptsoftware, die auf dem Pepper Roboter ausgeführt wird und ihn intern steuert. Diese kann mit persönlichen Modulen weiterentwickelt und angepasst werden. Dazu können mit Hilfe der Aldebaran \ac{SDK} auch NQOqi-\ac{SDK} genannt, eigene Module oder Bibliotheken entwickelt werden. Diese NAOqi-\ac{SDK} ist die Basis des NAOqi Frameworks und ist in C++ implementiert\cite{nao_dev_install_guide}.\\
Das NAOqi Framework ist das Programmiergerüst, welches zur Programmierung von NAO und Pepper Robotern verwendet wird. Es implementiert alle allgemeinen Anforderungen der Robotik, einschließlich: Parallelität, Ressourcen-Management, Synchronisation und Ereignisse. Dieses Framework ermöglicht eine homogene Kommunikation zwischen verschiedenen Modulen wie etwa die Bewegung, Audio oder Video sowie eine homogene Programmierung und einen homogenen Informationsaustausch. Das Framework ist:\\
\begin{itemize}
    \item plattformübergreifend, wie bereits erwähnt, basiert die NAOqi-\ac{SDK} auf C++ und bietet damit die Möglichkeit auf Windows, Linux oder sogar auch Mac zu entwickeln.
    \item sprachübergreifend, mit einer identischen API für C++ und Python. Weitere Details dazu sind in \autoref{subsubsec:Sprachuebergreifend} aufgeführt.
    \item fähig auf Introspektion, was bedeutet, dass das Framework weiß, welche Funktionen in den verschiedenen Modulen verfügbar sind und wo. Für Details diesbezüglich siehe \autoref{subsubsec:Introspektion}.
\end{itemize}

\subsubsection{Sprachübergreifend}\label{subsubsec:Sprachuebergreifend}
Software kann in C++ und Python entwickelt werden. Eine Übersicht über die Sprachen selbst in den Abschnitten \autoref{subsec:Cpp} und \autoref{subsec:Python}. In allen Fällen sind die Programmiermethoden genau die gleichen, alle vorhandenen \acp{API} können unabhängig von den unterstützten Sprachen aufgerufen werden:\\
\begin{itemize}
    \item Wird ein neues C++-Modul erstellt, können die C++-\ac{API}-Funktionen von überall aus aufgerufen werden,
    \item Sind sie richtig definiert, können auch die \ac{API}-Funktionen eines Python-Moduls von überall aus aufgerufen werden.
\end{itemize}
Normalerweise werden die Verhaltensweisen in Python und Ihre Dienste in C++ entwickelt.\\
\normfig{NAOqi-Framework-Sprachübergreifend}{0.5}{-}

\subsubsection{Introspektion}\label{subsubsec:Introspektion}
Die Introspektion ist die Grundlage der Roboter-\ac{API}, der Fähigkeiten, der Überwachung und der Maßnahmen bei überwachten Funktionen. Der Roboter selbst kennt alle verfügbaren \ac{API}-Funktionen. Wird eine Bibliothek entladen, werden die entsprechenden \ac{API}-Funktionen automatisch ebenfalls entfernt. Eine in einem Modul definierte Funktion kann der \ac{API} mit einem \texttt{BIND\_METHOD} hinzugefügt werden.\\
Wird eine Funktion gebunden, werder automatisch folgende Funktionen ausgeführt:\\
\begin{itemize}
    \item Funktionsaufruf in C++ und Python, wie in \autoref{subsubsec:Sprachuebergreifend} beschrieben
    \item Erkennen der Funktion, wenn sie gerade ausgeführt wird
    \item Funktion lokal oder aus der Ferne, z.B. von einem Computer oder einem anderen Roboter, ausführen weiter im Detail beschrieben in \autoref{subsubsec:VerteilterBaumUndKommunikation}
    \item Generierung und Aufruf von \texttt{wait}, \texttt{stop}, \texttt{isRunning} in Funktionen
\end{itemize}
Die \ac{API} wird im Webbrowser angezeigt wenn auf das Gerät per \ac{URL} oder \ac{IP}-Addresse auf dem Port 9559 zugegriffen wird. In dieser Übersicht, zeigt der Roboter seine Modulliste, Methodenliste, Methodenparameter, Beschreibungen und Beispiele an. Der Browser zeigt auch parallele Methoden an, die überwacht, zum Warten veranlasst und gestoppt werden können.\\
Die Introspektion und derer Implementation im NAOqi-Framework, ist also ein nützliches Werkzeug, welches es ermöglicht, die Roboter-\ac{API} zu verstehen und zu verwenden aber auch zu überwachenund zu steuern.\\
\normfig{NAOqi-Framework-Introspektion}{0.5}{-}

\subsubsection{Verteilter Baum und Kommunikation}\label{subsubsec:VerteilterBaumUndKommunikation}
Eine Echtzeitanwendung kann aus einer einzelnen ausführbaren Datei oder einem Baum von mehreren Systemen wie etwa Robotern, Prozessen oder Modulen bestehen. Unabhängig davon sind die Aufrufmethoden immer dieselben. Eine ausführbare Datei kann durch eine Verbindung mit einem anderen Roboter mit \ac{IP}-Adresse und Port verbunden werden, sodass alle \ac{API}-Methoden von anderen ausführbaren Dateien sind auf die gleiche Weise verfügbar sind, genau wie bei einer lokalen Methode. NAOqi trifft dabei selbst die Wahl zwischen schnellem Direktaufruf \ac{LPC} und Fernaufruf \ac{RPC}.\\
\normfig{NAOqi-Framework-Kommunikation}{0.5}{-}

\subsection{NAOqi Vorgehensweise}\label{subsec:NAOqiVorgehensweise}
Die NAOqi Software, welche auf dem Roboter läuft, ist ein Broker. Wenn dieser startet, lädt er eine Voreinstellungsdatei in den Speicher, in der festgelegt ist, welche Bibliotheken in dieser Konfiguration geladen werden sollen. Jede Bibliothek enthält ein oder mehrere Module, die den Broker benutzen, um ihre Methoden bereitzustellen.\\
\normfig{NAOqi-Broker-Bibliotheken-Module}{0.5}{-}
Der Broker selbst bietet Nachschlagdienste an, so dass jedes Modul im Baum oder im Netzwerk jede Methode finden kann, die an dem Broker bekannt gegeben wurde.\\
Das Laden von Modulen bildet dann einen Baum von Methoden, die mit Modulen verbunden sind, und von Modulen, welche mit dem Broker verbunden sind.\\
\normfig{NAOqi-Broker-Modul-Methoden}{0.5}{-}

\subsubsection{NAOqi Proxy}\label{subsubsec:NAOqiProxy}
Ein weiterer wichtiger Bestandteil des Broker ist der Proxy welcher grundlegend die Aufgabe eines Modules repräsentiert. Dies kann lokal oder entfernt enstehen der einzige Unterschied dabei ist, dass bei einer entfernten Refernz die \ac{IP} des entfernten Broker mit angegeben werden muss, gleich bleibt aber im internen, dass der Proxy die Methoden des Moduls an den Broker weiterleitet.\\

\subsection{NAOqi Module}\label{subsubsec:NAOqiModule}

\section{Robot Operating System}\label{sec:ROS}

\section{Programming Languages}

\subsection{C++}\label{subsec:Cpp}

\subsection{Python}\label{subsec:Python}

\section{Virtual Reality}

\section{Entwicklung für Virtual Reality}
\section{TCP}
\chapter{Technologieauswahl}
\section{VR-Brillen}
Die Auswahl der Technologie für die Entwicklung von Virtual-Reality-Anwendungen ist von entscheidender Bedeutung für den Erfolg eines Projekts. Bei der DHBW wurde die Entscheidung getroffen, die Metaquest 3 für das VR-Projekt zu verwenden. Diese Entscheidung wurde aufgrund mehrerer Faktoren getroffen, die im Folgenden erläutert werden.
\\

\noindent
Die Metaquest 3 wurde aufgrund ihrer vielseitigen Anwendungsmöglichkeiten und ihrer Benutzerfreundlichkeit ausgewählt. Die DHBW legt großen Wert darauf, den Studierenden eine moderne und zugängliche Lernumgebung zu bieten. Die Metaquest 3 erfüllt diese Anforderungen durch ihre intuitive Bedienung und ihre Fähigkeit, komplexe VR-Erlebnisse bereitzustellen, ohne dass zusätzliche Hardware wie externe Sensoren erforderlich sind.
\\

\noindent
Ein weiterer wichtiger Faktor bei der Auswahl der Metaquest 3 war die Integration von Oculus in die bestehende IT-Infrastruktur der DHBW. Die Unterstützung und Zusammenarbeit mit Oculus ermöglichte es der DHBW, Schulungen und Support für die Verwendung der Metaquest 3 bereitzustellen. Dies erleichterte die Einführung der VR-Technologie in den Lehrplan und sorgte für eine nahtlose Integration in bestehende Lehr- und Lernaktivitäten.
\\

\noindent
Darüber hinaus bietet die Metaquest 3 eine breite Palette von Anwendungen und Inhalten über den Oculus Store, einschließlich Bildungs- und Trainingsanwendungen, die für den Einsatz in der Hochschulbildung geeignet sind. Die Verfügbarkeit von hochwertigen Bildungsressourcen spielte eine wichtige Rolle bei der Entscheidung für die Metaquest 3, da sie den Lehrern und Studierenden Zugang zu einer Vielzahl von Lernmaterialien und Simulationen bietet, die den Lernprozess unterstützen und verbessern können.
\\

\noindent
Insgesamt wurde die Metaquest 3 aufgrund ihrer Benutzerfreundlichkeit, Integration in die bestehende Infrastruktur der DHBW und der Verfügbarkeit von Bildungsressourcen als ideale Wahl für das VR-Projekt der DHBW angesehen. Diese Auswahl bietet nicht nur eine solide Grundlage für die Entwicklung von VR-Anwendungen, sondern ermöglicht es auch, die VR-Technologie effektiv in den Lehrplan zu integrieren und den Lernerfolg zu steigern.
\section{Entwicklung für VR-Brillen}
Die Entscheidung, Unity als Entwicklungsplattform für das VR-Projekt "Pepper VR – Teleoperation eines humanoiden Roboters auf Basis der Analyse menschlicher Bewegung" zu wählen, wurde nach sorgfältiger Abwägung verschiedener Faktoren getroffen, wobei insbesondere auch die Unreal Engine in Betracht gezogen wurde.

\subsubsection{Unity}
Im folgenden werden die Vor- und Nachteile von Unity beleuchtet.
\paragraph{Vorteile:}

\begin{enumerate}
\item \textbf{Branchenübliche Plattform}: Unity ist eine der führenden Entwicklungsplattformen für VR-Anwendungen und wird von einer großen Community von Entwicklern und Unternehmen weltweit genutzt. Diese weitverbreitete Akzeptanz macht Unity zu einer branchenüblichen Wahl für die Entwicklung von VR-Inhalten und bietet Zugang zu einer Vielzahl von Ressourcen, Tutorials und Support, die für die erfolgreiche Umsetzung des Projekts entscheidend sind.
\item \textbf{Umfangreiche Funktionalitäten}: Unity bietet eine umfangreiche Auswahl an Funktionen und Werkzeugen, die speziell für die Entwicklung von VR-Anwendungen konzipiert sind. Die Integration von VR-Technologien wie Oculus Rift, HTC Vive und Metaquest in Unity ermöglicht es den Entwicklern, immersive VR-Erlebnisse mit hoher Qualität zu erstellen. Darüber hinaus bietet Unity eine benutzerfreundliche Oberfläche und eine intuitive Entwicklungsumgebung, die auch für Anfänger leicht zugänglich ist.

\item \textbf{Plattformübergreifende Unterstützung}: Unity ermöglicht die Entwicklung von VR-Anwendungen, die auf einer Vielzahl von Plattformen ausgeführt werden können, einschließlich PC, Konsolen, Mobilgeräten und Webbrowsern. Diese Flexibilität eröffnet die Möglichkeit, das VR-Projekt auf verschiedenen Geräten und Betriebssystemen zu testen und bereitzustellen, um eine maximale Reichweite und Zugänglichkeit zu gewährleisten.

\item \textbf{Erweiterbarkeit und Anpassbarkeit}: Unity zeichnet sich durch seine Erweiterbarkeit und Anpassbarkeit aus. Durch den Einsatz von Plugins und Assets aus dem Unity Asset Store können Entwickler zusätzliche Funktionen und Ressourcen in ihre VR-Anwendungen integrieren, was die Entwicklung beschleunigt und die Produktivität erhöht. Darüber hinaus bietet Unity die Möglichkeit, eigene Tools und Skripte zu erstellen, um die spezifischen Anforderungen des Projekts zu erfüllen und maßgeschneiderte Lösungen zu entwickeln.
\end{enumerate}

\paragraph{Nachteile:}

\begin{enumerate}
\item \textbf{Grafische Qualität}: Obwohl Unity in Bezug auf die Grafikqualität fortschrittliche Techniken bietet, erreicht es möglicherweise nicht das gleiche grafische Niveau wie die Unreal Engine, insbesondere bei fotorealistischen Rendering-Anforderungen.
\item \textbf{Lernkurve}: Unity kann eine steilere Lernkurve haben als die Unreal Engine, insbesondere für Anfänger oder Personen ohne Programmiererfahrung. Die Vielzahl von Funktionen und Optionen kann anfangs überwältigend sein.
\end{enumerate}

\subsubsection{Unreal Engine}
Nun werden die Vor- und Nachteile von Unreal Engine betrachtet.
\paragraph{Vorteile:}

\begin{enumerate}
\item \textbf{Grafische Qualität}: Die Unreal Engine ist bekannt für ihre beeindruckende Grafikqualität und ihre Fähigkeit, fotorealistische Umgebungen zu erstellen. Sie bietet fortschrittliche Rendering-Techniken wie Raytracing und hochwertige Materialien, die für VR-Anwendungen mit hohen grafischen Anforderungen von Vorteil sind.
\item \textbf{Visuelle Skripting-Tools}: Die Unreal Engine bietet visuelle Skripting-Tools wie den Blueprint-Editor, die es auch Personen ohne umfangreiche Programmierkenntnisse ermöglichen, komplexe Logik und Interaktionen zu erstellen. Dies kann die Entwicklungszeit verkürzen und die Kreativität fördern.

\item \textbf{Leistung}: Die Unreal Engine ist für ihre hohe Leistung und Stabilität bekannt, insbesondere bei großen Projekten mit komplexen Szenen und großen Datenmengen.
\end{enumerate}

\paragraph{Nachteile:}

\begin{enumerate}
\item \textbf{Einschränktere Plattformunterstützung}: Im Vergleich zu Unity bietet die Unreal Engine möglicherweise eine eingeschränktere Plattformunterstützung für die Entwicklung von VR-Anwendungen. Die Unterstützung für bestimmte VR-Geräte oder Plattformen kann begrenzt sein.
\item \textbf{Komplexität}: Die Unreal Engine kann aufgrund ihrer fortschrittlichen Funktionen und der visuellen Komplexität ihrer Benutzeroberfläche für Anfänger schwieriger zu erlernen sein. Die Blueprint-Logik kann zwar visuell sein, erfordert jedoch immer noch ein Verständnis von Konzepten wie Variablen und Logik.
\end{enumerate}

\subsubsection{Entscheidung}

Trotz der Vorteile der Unreal Engine in Bezug auf Grafikqualität und Leistung wurde Unity als die bevorzugte Entwicklungsplattform für das VR-Projekt "Pepper VR – Teleoperation eines humanoiden Roboters auf Basis der Analyse menschlicher Bewegung" gewählt. Die breite Unterstützung, die umfangreichen Funktionalitäten, die plattformübergreifende Unterstützung und die Erweiterbarkeit von Unity waren entscheidend für diese Wahl. Unity bietet eine solide Grundlage für die Entwicklung hochwertiger VR-Anwendungen, die den Anforderungen des Projekts gerecht werden und eine erfolgreiche Integration von VR-Technologie in den Lehrplan ermöglichen.
\subsection{Verbindungstechnologie}
Die Wahl von Unity als Entwicklungsplattform für das VR-Projekt "Pepper VR – Teleoperation eines humanoiden Roboters auf Basis der Analyse menschlicher Bewegung" führte zur Notwendigkeit, eine Methode zur Kommunikation zwischen der Unity-Anwendung und dem humanoiden Roboter Pepper zu implementieren. Die ROS TCP-Verbindung wurde aufgrund ihrer Kompatibilität mit Unity und der Robotersteuerung über das Robot Operating System (ROS) als geeignete Lösung identifiziert.

\subsubsection{Kompatibilität mit Unity}

ROS TCP (Transmission Control Protocol) bietet eine zuverlässige Methode zur Datenübertragung zwischen ROS und Unity. Unity unterstützt die Kommunikation über TCP/IP-Sockets, was es ermöglicht, Daten zwischen der Unity-Anwendung und externen Geräten wie Robotern über das Netzwerk auszutauschen. Durch die Implementierung einer ROS TCP-Verbindung kann die Unity-Anwendung Befehle an den Roboter senden und Daten von ihm empfangen, was eine nahtlose Integration in das VR-Erlebnis ermöglicht.

\subsubsection{Robotersteuerung über ROS}

Pepper, der humanoide Roboter, wird über das Robot Operating System (ROS) gesteuert, das eine Standardplattform für die Entwicklung von Robotersoftware ist. ROS bietet eine Vielzahl von Funktionen zur Robotersteuerung, einschließlich der Unterstützung für verschiedene Sensoren, Aktuatoren und Navigationssysteme. Indem die ROS TCP-Verbindung verwendet wird, kann die Unity-Anwendung mit den ROS-Nodes kommunizieren, die für die Steuerung von Pepper zuständig sind. Dies ermöglicht es der Unity-Anwendung, Pepper-Bewegungen zu steuern und Sensordaten von Pepper zu empfangen, um ein interaktives VR-Erlebnis zu schaffen.

\subsubsection{Implementierung in Unity}

Die Implementierung der ROS TCP-Verbindung in Unity erfolgt mithilfe von Plugins oder eigenen Skripten, die die TCP/IP-Kommunikation ermöglichen. Durch die Verwendung von vorhandenen ROS-Bibliotheken oder der Erstellung benutzerdefinierter ROS-Nodes kann die Unity-Anwendung ROS-Nachrichten senden und empfangen, um mit Pepper zu interagieren. Dies ermöglicht es, die Bewegungen von Pepper in Echtzeit zu steuern und Feedbackdaten von Pepper in die Unity-Anwendung zu integrieren.
\\

\noindent
Insgesamt wurde die ROS TCP-Verbindung aufgrund ihrer Kompatibilität mit Unity und der Robotersteuerung über ROS als ideale Lösung für die Kommunikation zwischen der Unity-Anwendung und dem humanoiden Roboter Pepper ausgewählt. Diese Entscheidung ermöglicht es, eine immersive und interaktive VR-Erfahrung zu schaffen, bei der die Benutzer direkt mit dem Roboter interagieren können.
\chapter{Umsetzung}
\section{MetaQuest3}

\section{ROS-Topics}\label{sec:ros-topics}
Die allgemeine Datenkommunikation des Projektes wird zwischen den Komponenten ausschließlich über \ac{ROS}-Topics abgewickelt und über \ac{ROS} verwaltet. Welche dabei standardmäßig von welchen Komponenten angeboten werden und welche noch zusätzlich hinzugefügt werden müssen werden im folgenden Abschnitt beschrieben.\\
\subsection{Videostream und Kamerabilder}\label{subsec:Videostream}
(nicht vollständig Implementiert)\\\\
Einen Videostream wie man ihn von einem normalen Video-Player kennt, ist in \ac{ROS} nicht direkt möglich. Stattdessen wird der Stream als eine Abfolge von Bildern übertragen. Diese Bilder werden in einer Node, die Zugriff auf die Kamera hat, aufgenommen, nach Bedarf verarbeitet und schließlich über ein \ac{ROS}-Topic veröffentlicht.\\
Dieser pseudo-Stream wird standardmäßig von der \textit{NAOqi\_bridge} bereitgestellt. Dadurch ist nicht nur der Zugriff auf die front-Kamera des Roboters gegeben, sondern alle Kameras und Sensoren des Roboters. Bei den Kameras kann zwischen der front-Kamera und der bottom-Kamera gewählt werden, wobei die NAOqi\_bridge die Bilder der Kameras in verschiedenen Topics veröffentlicht wodurch verschiedene Features der Kameras genutzt werden können.\\
Die verschiedenen Topics sind kaskadiert aufgebaut und liegen unter dem Topic \textit{/pepper\_robot/camera/front/}, \textit{/pepper\_robot/camera/bottom/} und \textit{/naoqi\_bridge/camera/depth}  wobei für das Projekt nur die Frontkamera relevant ist. Sie bietet folgende Topics an:
\begin{itemize}
    \item \textit{image\_raw} - Das unverarbeitete Bild der Kamera in voller Auflösung und Farbtiefe
    \item \textit{image\_raw/compressed} - Das komprimierte Bild der Kamera komprimiert durch die \ac{JPEG}-Kompression
    \item \textit{image\_raw/theora} - Das komprimierte Bild der Kamera im Theora-Format ebenfalls komprimiert mit Theora
\end{itemize}
Das Projekt nutzt das Topic \textit{image\_raw/compressed} um die Bilder der Kamera zu empfangen. Eine direkte Übertragung der Bilder unkomprimiert, ist in der Theorie zwar möglich ist aber für eine \ac{AR} Anwendung nicht von Vorteil, da die Bilder in der Regel nicht in voller Auflösung benötigt werden und die Übertragung der Bilder unkomprimiert zu einer hohen Netzwerkauslastung führen würde, was wiederum die Latenz der Bilder erhöhen und die Anzahl der Bilder pro Sekunde begrenzen würde. Dies wiederum würde im Umkehrschluss die Qualität der \ac{AR} Anwendung drastisch verschlechtern und ein magelhaftes Erlebnis bieten, welches bis hin zur Übelkeit des Nutzers reichen könnte.\\
Die Bilder werden von der NAOqi\_bridge in einem festen Intervall von 30 Bildern pro Sekunde veröffentlicht. Diese Rate ist bekannt als die niedrigste Rate, die der Mensch als flüssige Bewegung wahrnimmt und damit für die Anwendung ausreichend.\\
Mit einer \ac{HD} Auflösung von 1280x720 Pixel, 8-Bit Farbtiefe und einer angenommenen Kompressionsrate von 0.5 durch \ac{JPEG}, ergibt sich bei den genannten 30 Bildern pro Sekunde unkomprimiert eine Datenrate von etw. 663Mbit/s und mit Kompression etwa 331.5Mbit/s.\\
\begin{equation}
    1280px \cdot 720px \cdot 8bit \cdot 30fps = 663Mbit/s
    \label{eq:datenrate}
\end{equation}
\begin{equation}
    633Mbit/s \cdot 0.5 = 331.5Mbit/s
    \label{eq:datenrate-komprimiert}
\end{equation}

\subsection{Teleoperationspositionen}
(nicht vollständig Implementiert)\\\\
Der Roboter selbst besitzt drei maßgebende Kinematiken, welche für die Teleoperation genutzt werden. Diese sind die beiden Arme  \textit{LArm} und \textit{RArm} und der Kopf \textit{Head}. Eine Ausnahme bildet das Fahrwerk, welches nicht direkt über Teleoperation gesteruert wir, aber trotzdem eine unabdingliche Rolle für die Teleoperation und dem Nutzen des Roboters spielt. Über die NAOqi\_bridge können Positionen der einzelnen Gelenke der Kinematiken abgefragt und gesetzt werden, was für das Fahrwerk und das Kopfgelenk direkt genutzt werdenk kann, jedoch für die Arme aber durch die hohe Anzahl an Gelenken zuerst weiter verarbeitet werden muss.\\
Die Rohdaten für die gewünschte Position der Kinematiken werden von der MetaQuest3 Anwendung beziehungsweise Node bereitgestellt und über das Topic \textit{/AR/position/} mit den sub-Topics \textit{/LController}, \textit{/RController} und \textit{/Head} veröffentlicht. Wiederum eine Ausnahme bildet das Fahrwerk, welches indirekt seine Daten aus den Kontroller-Knöpfen der MetaQuest3 bezieht und unter dem Topic \textit{/AR/button} veröffentlicht.\\
Eine Position besteht sowohl aus der aktuellen und relativen Position der Komponente zu der Brille, als auch aus der aktuellen Winkel der Komponente. Diese beiden Werte werden in der Anwendung in einem festen Format gespeichert und über die Topics, einzeln für jede Kinematik-Gruppe veröffentlicht. Die dazugehörige Nachricht ist in Listing \autoref{lst:position-msg} zu sehen.\\
\begin{lstlisting}[language=Python, caption=Positions Topic, label=lst:position-msg]
    float64 x
    float64 y
    float64 z
    float64 roll
    float64 pitch
    float64 yaw
\end{lstlisting}
Im Gegensatz dazu die Winkel um die Achsen x, y und z, zu benennen ist es in der Robotik üblich diese als Roll, Pitch und Yaw zu bezeichnen. Diese Werte sind daher auch in der Nachricht als \textit{roll}, \textit{pitch} und \textit{yaw} zu finden. Sie werden statt in Winkel 0-360Grad in Radiant angegeben und können von -$\pi$ bis $\pi$ reichen. Die Positionen werden ebenfalls gemäß dem Standard in Metern angegeben und können in dem im Headset genullten Koordinatensystem frei bewegt werden.\\
Wieder die Ausnahme bildet das Fahrwerk, welches nicht in einem 3D Raum bewegt wird, sondern nur in einer Ebene. Daher wird die Position des Fahrwerks nur in x und y angegeben und die Winkel nur in yaw. Die restlichen Werte werden ignoriert und für die Berechnung auf 0 gesetzt.\\

\section{Streaming-Node}\label{sec:Streaming-Node}
(nicht vollständig Implementiert)\\\\
Für die Übertragung der Kamerabilder wie sie in \autoref{subsec:Videostream} beschrieben sind, wird eine eigene Node für die MetaQuest3 erstellt. Sie hat die einzig und alleinige Aufgabe die Kamerabilder des Pepper, aus der, von der NAOqi\_bridge zur Verfügung gestellten Topic \\ \textit{/pepper\_robot/camera/front/compressed} aufzunehmen, in einem möglichst kleinen Zwischenspeicher (Buffer) zu speichern und mit möglichst geringer Latenz an die Brille weiter zu senden.\\
Dazu wird in der \ac{AR}-App (die \ac{ROS} Node), die Topic \textit{/pepper\_robot/camera/front/image\_raw/compressed} abonniert. Die Bilder werden von dieser Topic so schnell wie es \ac{ROS} ermöglicht abgegriffen, und ind den Buffer geschrieben. Dadurch kann es zwar vorkommen, dass das selbe Bild mehrere male aus der Topic gelesen wird, jedoch ist es eine einfache Methode die Bilder zu übertragen und die Latenz gering zu halten um wiederum den pseudo-Stream so flüssig und responsiv zu gestalten wie es nur möglich ist. Dadurch ist der Limitierende Faktor das darunter liegende System oder primär die Bandbreite der Verbindung im Netzwerk.\\
Der eben erwähnte Buffer ist nicht zwingend notwendig, aber für die Anwendung von Vorteil, da die Bilder so nicht direkt an das Display weitergegeben werden müssen, sondern in einem Zwischenspeicher gespeichert werden können. Dadurch können die kleinen Schwankungen in der Übertragungsgeschwindigkeit der Bilder ausgeglichen werden und die Bilder können so gleichmäßig wie möglich angezeigt werden. Jedoch bedeutet ein Buffer auch, dass die Bilder nicht in Echtzeit angezeigt werden können, sondern immer mit einer kleinen Verzögerung, weswegen der Buffer sebst so klein wie möglich gehalten werden sollte. Daher wurde sich in diesem Projekt für einen Ring-Buffer entschieden, welcher lediglich die letzten 5 Bilder speichert, da eine Verzögerung von 5 Bildern (bei 30\ac{FPS} etwa 0.166s) zwar schon als unnatürlich und störend wahrgenommen werden kann, aber in diesem Aufbau in Kauf genommen wird. Auf Verzögerungen durch die Trägheit der Kinematiken des Roboters wird in im späterne Verlauf unter \autoref{sec:Berechnungs-Node} nocheinmal separat eingegangen.\\
Eine Besonderheit der Node ist, dass Sie nicht klassisch unter \textit{rclcpp} oder \textit{rclpy} implementiert wird, sondern über den \ac{ROS}-\ac{TCP}-Connector. Dieser ist ein Modul, bereitgestellt durch Unity, welches die Kommunikation zwischen \ac{ROS} und Unity ermöglicht. Dadruch werden die Nachrichten der \ac{ROS}-Topics direkt in Unity empfangen und können dort unter der Verwendung von C\# weiterverarbeitet werden. Durch die Nutzung der Programiersprache C\# verläuft das Empfangen und Speichern der Bilder in Unity deutlich schneller und effizienter als in Python, was wiederum die Latenzen der Bilder reduziert und die Anwendung flüssiger macht.\\
Weiter wird ein Modul mit dem Namen \textit{openXR} genutzt, welches die Kommunikation zwischen Unity und der MetaQuest3 erheblich erleichter und standartisiert. Weiter ermöglicht es die App auf einem Windows-Rechner zu entwickeln und auszuführen, sodass die App nicht auf der Brille selbst ausgeführt wird. Dadurch kann Rechenleistung und Speicherplatz auf der Brille gespaart werden und die App kann auf einem leistungsstärkeren Rechner ausgeführt werden.\\

Für die Erstellung der App wird ein neues Unity-Projekt erstellt und die benötigten Module \textit{openXR} und \textit{ROS-TCP-Connector} hinzugefügt. Die Module werden in Unity als \textit{.unitypackage} Datei bereitgestellt und werden über den Menüpunkt \textit{Assets} $\rightarrow$ \textit{Import Package} $\rightarrow$ \textit{Custom Package} hinzugefügt werden.\\
Der Kern der Unity-App besteht aus der in \autoref{lst:unity-controller-data} gezeigten Klasse. Diese hier beispielhaft für den rechten Controller gezeigte Klasse, ist für die Aufnahme der Daten des Controllers zuständig und veröffentlicht diese in einem \ac{ROS}-Topic. Die Klasse wird in Unity als \textit{MonoBehaviour} implementiert und über die Funktionen \textit{Start()} und \textit{Update()} gesteuert. Die Funktion \textit{Start()} wird beim Start der App aufgerufen und initialisiert die Verbindung zu \ac{ROS} und registriert den Publisher für die Nachrichten. Die Funktion \textit{Update()} wird in jedem Frame aufgerufen und überprüft ob die Nachrichten veröffentlicht werden sollen. Die Nachrichten werden dabei in einem festen Intervall von 0.1s veröffentlicht, was wiederum der maximalen Frequenz des Peppers entspricht, und somit eine flüssige und efiziente Bewegung ermöglicht.\\
\begin{lstlisting}[language=csh, caption=Abgreifen der Controller Daten, label=lst:unity-controller-data]
    using UnityEngine;
    using Unity.Robotics.ROSTCPConnector;
    using RosMessageTypes.Geometry;
    
    public class ControllerPositionPublisher : MonoBehaviour
    {
        ROSConnection ros;
        public string RController = "/AR/RController";
        public GameObject rightController;
        
        public float publishMessageFrequency = 0.1f;
    
        void Start()
        {
            ros = ROSConnection.GetOrCreateInstance();
            ros.RegisterPublisher<PoseStampedMsg>(RController);
            timeElapsed = 0;
        }
    
        void Update()
        {
            timeElapsed += Time.deltaTime;
    
            if (timeElapsed > publishMessageFrequency)
            {
                ros.Publish(topicName,
                    GetControllerPose(rightController));
    
                timeElapsed = 0;
            }
        }
    
        PoseStampedMsg GetControllerPose(GameObject controller)
        {
            PoseStampedMsg poseMsg = new PoseStampedMsg();
            poseMsg.header.frame_id = "unity";
            poseMsg.header.stamp
                = ROSConnection.GetOrCreateInstance()
                    .GetROSTime();
            
            poseMsg.pose.position = new PointMsg(
                controller.transform.position.x,
                controller.transform.position.y,
                controller.transform.position.z
            );
            
            poseMsg.pose.orientation = new QuaternionMsg(
                controller.transform.rotation.x,
                controller.transform.rotation.y,
                controller.transform.rotation.z,
                controller.transform.rotation.w
            );
    
            return poseMsg;
        }
    }
\end{lstlisting}
Das Bild der Kamera ist dagegen hat zu der einfachen Struktur des publishen der Controller-Daten eine weitere Komponente, den Ring-Buffer. Auch hier werden wieder die üblichen \ac{ROS} Funktionen initialiert. In der \textit{Start()} Funktion wird der Publisher für die Kamerabilder (z.9) registriert und der Buffer, wie in \autoref{lst:unity-image-data}(z.24) mit kleinen temporären Texturen gefüllt, initialisiert. Dazu werden Behelfsvariablen für die Größe des Buffers und die aktuelle Position im Buffer erstellt. In der \textit{Update()} Funktion wird dann in jedem Frame das Bild der Kamera in den Buffer geschrieben und die Nachricht veröffentlicht. Sobald in der Topic der Kamera ein Update stattfindet, wird die Funktion \textit{UpdateImage()} aufgerufen und das Bild aus der Nachricht in den Buffer geschrieben (z.34) und dessen Index erhöht (z.35). Zum schluss der Methode wird das Bild aus dem Buffer in die Textur des Renderers geschrieben und angezeigt (z.37).\\
\begin{lstlisting}[language=csh, caption=Abgreifen der Kamerabilder, label=lst:unity-image-data]
    using UnityEngine;
    using Unity.Robotics.ROSTCPConnector;
    using RosMessageTypes.Sensor;
    using System;
    using System.Runtime.InteropServices;

    public class ImageSubscriber : MonoBehaviour
    {
        ROSConnection ros;
        public string topicName
            = "/pepper_robot/camera/front/image_raw/compressed";
        public Renderer imageRenderer;

        private Texture2D[] imageBuffer;
        private int bufferSize = 5;
        private int currentIndex = 0;

        void Start()
        {
            ros = ROSConnection.GetOrCreateInstance();
            ros.Subscribe<ImageMsg>(topicName, UpdateImage);

            imageBuffer = new Texture2D[bufferSize];
            for (int i = 0; i < bufferSize; i++)
            {
                imageBuffer[i]
                    = new Texture2D(2,
                        2,
                        TextureFormat.RGB24,
                        false);
            }
        }

        void UpdateImage(ImageMsg imageMessage)
        {
            Texture2D texture
                = new Texture2D(imageMessage.width,
                    imageMessage.height,
                    TextureFormat.RGB24,
                    false);
            texture.LoadRawTextureData(imageMessage.data);
            texture.Apply();

            imageBuffer[currentIndex] = texture;
            currentIndex = (currentIndex + 1) % bufferSize;

            imageRenderer.material.mainTexture
                = imageBuffer[currentIndex];
        }
    }
\end{lstlisting}


\section{Berechnungs-Node}\label{sec:Berechnungs-Node}
(nicht vollständig Implementiert)\\\\
Die Berechnung der Kinematiken und der genauen Positionen der einzelnen Motoren werden über inverse Kinematiken gelöst, welche von der \textit{pyBullet} Bibliothek bereitgestellt werden. Dazu wird einen individuelle Node erstellt, welche die \textit{pyBullet} Bibliothek nutzt und die berechneten Positionen an die NAOqi\_bridge weitergibt. Die Node hat dabei eine Art Middleware-Funktion und übernimmt die Umrechnungen zwischen der MetaQuest3 und der NAOqi\_bridge.\\
Die Node bindet die \textit{pyBullet} und die \textit{NAOqi\_bridge} Bibliotheken ein. Für die Berechnung abonniert sie als Input die Topics der MetaQuest3: \textit{/AR/position/LController}, \textit{/AR/position/RController}, \textit{/AR/position/Head} und \textit{/AR/button}. Der Output an die NAOqi\_bridge ist dagegen etwas breiter gefächtert, und ist eine direkte Kommunikation an die \ac{API}.\\

Die Node besteht nicht wie jede Node aus einem Publisher-Teil und einem Subscriber-Teil. Sondern nur dem Subscriber-Teil, da die Node nur die Positionen der MetaQuest3 empfängt und an die NAOqi\_bridge weitergibt. Die Berechnung der Positionen und die Kommunikation mit der NAOqi\_bridge erfolgt in einem einzigen Teil der Node. Dazu ist maßgebend die Funktion \textit{callback} von Relevanz, da diese Aufgerufen wird sobald eine Nachricht in einem der abonnierten Topics ankommt. Als zweite zentrale Komponente dient hier die \textit{pyBullet} Bibliothek. Sie muss zunächst initialisiert werden und wie in \autoref{lst:berechnungs-node-pybullet} gezeigt die Kinematiken des Roboters über die \ac{URDF}-Datei des Peppers einlesen.\\
\begin{lstlisting}[language=Python, caption=Berechnungs-Node-Initialisierung von PyBullet, label=lst:berechnungs-node-pybullet]
    import pybullet as p

    p.connect(p.DIRECT)

    robot = p.loadURDF("pepper.urdf",
            [0, 0, 0],
            useFixedBase=True)
\end{lstlisting}
Sobald eine Nachricht empfangen wird, werden die Nachrichten aus den Arrays der Nachricht der MetaQuest3 entnommen und durch nach der Berechnung der inversen Kinematiken in die passenden Werte für die NAOqi\_bridge umgerechnet. Das Array der Positionen der MetaQuest3 wird wie in \autoref{lst:position-msg} beschrieben aufgebaut und enthält die Werte für x, y, z, roll, pitch und yaw, in genau dieser Reihenfolge. Die Werte werden dabei direkt per \texttt{rospy.Subscriber(``[Topic]''), [Type], [callback]} abonniert und in der Funktion \textit{callback} verarbeitet. Hier werden für jede Gelenkgruppe (Arme und Genick) eine separate callback-Funktion aufgerufen. Diese Gruppen werden zu Begin in einem Array aus den einzelnen Gelenken zusammengefasst und der jeweiligen Funktion übergeben. Aus der \ac{URDF}-Datei des Pepper werden die Gelenke und deren Positionen ausgelesen (\autoref{lst:pepper-urdf}) und in einem Dictionary gespeichert. Die Positionen der Gelenke werden dann in einem weiteren Dictionary gespeichert und an die \textit{pyBullet} Bibliothek übergeben. Aus der \ac{URDF}-Datei des Pepper können die Namen der Gelenke gelesen werden, und die Gelenke können über die Namen angesprochen werden und den Gruppen zugeordnet werden. Dazu wird der Code aus \autoref{lst:pepper-urdf} und \autoref{lst:gelenk-grupperierung} genutzt.\\
\begin{lstlisting}[language=Python, caption=Ausschnitt der URDF-Datei des Pepper, label=lst:pepper-urdf]
    <joint name="HeadYaw" type="continuous">
        <parent link="base_link"/>
        <child link="HeadYaw_link"/>
        <origin xyz="0 0 0.126" rpy="0 0 0"/>
        <axis xyz="0 0 1"/>
        <limit lower="-2.0857" upper="2.0857"/>
    </joint>
\end{lstlisting}

\begin{lstlisting}[language=Python, caption=Gelenk-Grupperierung, label=lst:gelenk-grupperierung]
    joints = {
        "Head": ["HeadYaw", "HeadPitch"],
        "LArm": ["LShoulderPitch",
                "LShoulderRoll",
                "LElbowYaw",
                "LElbowRoll",
                "LWristYaw"],
        "RArm": ["RShoulderPitch",
                "RShoulderRoll",
                "RElbowYaw",
                "RElbowRoll",
                "RWristYaw"]
    }
\end{lstlisting}
Damit kann über \texttt{joints[``Gruppe'']} auf die einzelnen Gelenke zugegriffen werden und die Positionen der Gelenke gelesen und gesetzt werden. Die Berechnung der inversen Kinematiken erfolgt über die \textit{pyBullet} Bibliothek und die Funktion \texttt{p.calculateInverseKinematics([Robot], [Gelenk-Gruppe], [Position], [Orientierung])}. Dabei können die Positionen und Orientierungen der Gelenke direkt in der Funktion übergeben werden und die Funktion gibt die berechneten Positionen der Gelenke zurück.\\
Diese Daten sind nun so aufbereitet, dass sie direkt an die NAOqi\_bridge übergeben werden können. Dazu wird zunächst der eine Session des Roboters geöffnet und dieser der Service \textit{ALMotion} über \texttt{m\_service = session.service(ALMotion)} gestartet. Die Positionen der Gelenke werden dann über den Service und der Funktion \texttt{m\_service.setAngles([Name], [Winkel], [MaxGeschwindigkeit])} an die NAOqi\_bridge übergeben und die Gelenke des Roboters werden in die gewünschte Position gebracht. Unter den Namen werden in diesem Projekt die folgenden Gelenke angesprochen:
\begin{itemize}
    \item \textit{/pepper\_robot/HeadYaw} - Die Position des Kopfes um die Y-Achse.
    \item \textit{/pepper\_robot/HeadPitch} - Die Position des Kopfes um die X-Achse.
    \item \textit{/pepper\_robot/[R,L]ShoulderPitch} - Die Position des Schultergelenks um die X-Achse.
    \item \textit{/pepper\_robot/[R,L]ShoulderRoll} - Die Position des Schultergelenks um die Y-Achse.
    \item \textit{/pepper\_robot/[R,L]ElbowYaw} - Die Position des Ellenbogengelenks um die Y-Achse.
    \item \textit{/pepper\_robot/[R,L]ElbowRoll} - Die Position des Ellenbogengelenks um die X-Achse.
    \item \textit{/pepper\_robot/[R,L]WristYaw} - Die Position des Handgelenks um die Y-Achse.
    \item \textit{/pepper\_robot/WheelFL} - Die Position des linken Vorderrades.
    \item \textit{/pepper\_robot/WheelFR} - Die Position des rechten Vorderrades.
    \item \textit{/pepper\_robot/WheelB} - Die Position des hinteren Rades.
\end{itemize}

Diese Daten werden mit der maximalen Frequenz des Peppers von 10Hz an die NAOqi\_bridge übergeben womit bedingt durch die Bewegung des Users und die Trägheit eine Art Interpolation geschaffen wird und die Bewegungen des Roboters flüssiger und natürlicher wirken.\\

\section{Fahrwerk-Node}\label{sec:Fahrwerk-Node}
(nicht vollständig Implementiert)\\\\
Das Fahrwerk des Pepper wird nicht direkt über die MetaQuest3 gesteuert, sondern über die Knöpfe der MetaQuest3. Dadurch kann hier weder mit Positionen gearbeitet werden, noch mit inversen Kinematiken. Ebenfalls platziert in der Berechnungs-Node, können die Variablen wie die \texttt{session} und der \texttt{m\_service} auch für diese Berechnungen genutzt werden. Für diesen Ansatz wird jedoch eine separate callback-Funktion benötigt. Sie reagiert auf die \textit{/AR/button} Nachrichten und lässt den Roboter entsprechend der Knöpfe bewegen. Hierzu werden keine aufwendigen Berechnungen benötigt, sondern nur die Knöpfe ausgelesen und eine Translation oder Rotation des Roboters durchgeführt.\\
Die Knöpfe der MetaQuest3 sind in diesem Projekt wie folgt belegt:
\begin{itemize}
    \item \textit{A} - Vorwärtsbewegung
    \item \textit{B} - Rückwärtsbewegung
    \item \textit{X} - Drehung nach links
    \item \textit{Y} - Drehung nach rechts
\end{itemize}
Auf die Knöpfe wird über die \textit{OVRInput}-Klasse von Unity zugegriffen und als einfacher charakter in der Topic übergeben. Dadurch kann zwar nur ein einzelner Befehl ausgeführt werden, was jedoch über die Tastenbelegung eher durch den User begrenzt wird. Auch das Fahrwerk kann über den \texttt{m\_service} angesprochen werden und die Räder des Roboters können über die Funktion \texttt{m\_service.moveToward([X], [Y], [Theta])} in die gewünschte Richtung bewegt werden. Die Werte für X und Y sind dabei die Geschwindigkeiten der Vorderräder und der Wert für Theta ist die Geschwindigkeit des hinteren Rades. Die Werte können dabei zwischen -1 und 1 liegen, wobei -1 die maximale Geschwindigkeit in die eine Richtung und 1 die maximale Geschwindigkeit in die andere Richtung bedeutet. Auch hier wird die Funktion mit 10Hz aufgerufen und der Roboter bewegt sich nur solange einer der Knöpfe gedrückt wird. Wird ein Knopf entlastet stoppt der Roboter. Die Bewegung des Roboters wird dabei durch die Trägheit des Roboters und die Reaktionszeit des Users begrenzt und wirkt so natürlicher und flüssiger.\\

\chapter{Anwendungsgebiete}

Die Teleoperation humanoider Roboter bietet eine Vielzahl von Einsatzmöglichkeiten in unterschiedlichen Bereichen. Im Folgenden werden einige der wichtigsten Anwendungsgebiete beschrieben:

\subsection{Medizin und Pflege}
In der Medizin und Pflege können humanoide Roboter zur Unterstützung von Ärzten und Pflegepersonal eingesetzt werden. Sie können in gefährlichen oder infektiösen Umgebungen arbeiten, wodurch das Risiko für medizinisches Personal minimiert wird. Zudem können Roboter Routineaufgaben übernehmen, wie beispielsweise die Ausgabe von Medikamenten oder das Überwachen von Vitalparametern, was die Arbeitsbelastung des Personals verringert und mehr Zeit für die direkte Patientenbetreuung ermöglicht.

\subsection{Industrielle Fertigung}
In der industriellen Fertigung können humanoide Roboter für Aufgaben eingesetzt werden, die für Menschen gefährlich oder ergonomisch ungünstig sind. Sie können in gefährlichen Umgebungen arbeiten, schwere Lasten heben oder repetitive Aufgaben mit hoher Präzision ausführen. Durch die Teleoperation können Experten aus der Ferne eingreifen und die Roboter steuern, was die Flexibilität und Effizienz der Fertigungsprozesse erhöht. Darüber hinaus können sie in der Automobilindustrie für die Montage von Fahrzeugteilen eingesetzt werden, wo Präzision und Geschwindigkeit entscheidend sind.

\subsection{Katastrophenhilfe und Rettungseinsätze}
Humanoide Roboter können in Katastrophengebieten eingesetzt werden, um Menschen zu retten oder gefährliche Situationen zu erkunden. Sie können durch Trümmer navigieren, Verletzte finden und erste Hilfe leisten. Durch die Teleoperation können Rettungskräfte aus sicherer Entfernung arbeiten und dennoch präzise und effektiv handeln. Diese Roboter können auch bei der Suche nach Verschütteten nach Naturkatastrophen wie Erdbeben oder Erdrutschen eingesetzt werden, wo der Zugang für Menschen schwierig oder gefährlich ist.

\subsection{Raumfahrt}
In der Raumfahrt können humanoide Roboter für Außenbordeinsätze (EVAs) und Wartungsarbeiten an Raumstationen oder Satelliten eingesetzt werden. Sie können in extremen Umgebungen arbeiten, die für Menschen gefährlich sind, und komplexe Aufgaben mit hoher Präzision ausführen. Die Teleoperation ermöglicht es, dass die Roboter von der Erde aus gesteuert werden, wodurch das Risiko für Astronauten minimiert wird. Diese Roboter können auch bei der Exploration anderer Planeten eingesetzt werden, um Proben zu sammeln und wissenschaftliche Experimente durchzuführen.

\subsection{Bildung und Forschung}
Humanoide Roboter können in Bildungseinrichtungen und Forschungslabors eingesetzt werden, um Experimente durchzuführen und komplexe Konzepte zu demonstrieren. Sie können als Lehrassistenten dienen und Schülern und Studenten interaktive und praxisnahe Lernerfahrungen bieten. In der Forschung können Roboter zur Untersuchung neuer Technologien und zur Entwicklung innovativer Anwendungen verwendet werden. Durch die Teleoperation können Forscher und Studenten Roboter aus der Ferne steuern und Experimente in verschiedenen Umgebungen durchführen.

\subsection{Heim- und Unterhaltungselektronik}
Im Bereich der Heim- und Unterhaltungselektronik können humanoide Roboter als persönliche Assistenten oder Unterhaltungsgeräte eingesetzt werden. Sie können Aufgaben im Haushalt übernehmen, wie das Aufräumen oder das Zubereiten von Mahlzeiten, und bieten interaktive Unterhaltungsmöglichkeiten, wie das Spielen von Spielen oder das Führen von Gesprächen. Durch die Teleoperation können Benutzer ihre Roboter aus der Ferne steuern und ihnen Anweisungen geben, was die Interaktion mit dem Roboter und die Nutzung seiner Funktionen erleichtert.

\subsection{Kundendienst und Gastgewerbe}
Humanoide Roboter können im Kundendienst und im Gastgewerbe eingesetzt werden, um Kunden zu begrüßen, Informationen bereitzustellen oder Bestellungen entgegenzunehmen. Sie können in Hotels, Restaurants oder Geschäften arbeiten und den Service verbessern, indem sie rund um die Uhr verfügbar sind und in verschiedenen Sprachen kommunizieren können. Die Teleoperation ermöglicht es, dass Menschen aus der Ferne eingreifen und bei Bedarf Unterstützung leisten können. Dies erhöht die Effizienz des Kundendienstes und trägt zur Zufriedenheit der Kunden bei.
\\

\noindent
Die Anwendungsgebiete der Teleoperation humanoider Roboter sind vielfältig und bieten zahlreiche Möglichkeiten, die Effizienz und Sicherheit in verschiedenen Bereichen zu erhöhen. Die vorliegende Arbeit trägt dazu bei, diese Vision durch die Entwicklung einer VR-basierten Steuerung für den humanoiden Roboter Pepper weiter voranzutreiben. Durch die Integration von Virtual Reality und Robotik können neue Lösungen für komplexe Herausforderungen entwickelt werden, die das Potenzial haben, das Leben von Menschen weltweit zu verbessern.

\chapter{Fazit}

Die Umsetzung des Projekts \textit{Pepper VR – Teleoperation eines humanoiden Roboters auf Basis der Analyse menschlicher Bewegung} stellte eine anspruchsvolle Herausforderung dar, die tiefgreifende Kenntnisse in verschiedenen Bereichen der Informatik und Robotik erforderte. Trotz der ambitionierten Ziele, die Verbindung zwischen der MetaQuest 3 VR-Brille und dem humanoiden Roboter Pepper herzustellen und die Steuerung von Pepper durch die VR-Controller zu ermöglichen, stießen wir auf mehrere technische und organisatorische Hürden, welche die Umsetzung leider schlussendlich nicht möglich machten.

Die wichtigsten Erkenntnissen und Schlussfolgerungen sind:

\section{Technische Herausforderungen}
Die Implementierung der Steuerung eines humanoiden Roboters über eine VR-Brille wie die MetaQuest 3 erfordert eine präzise und latenzfreie Datenübertragung. Probleme mit der Netzwerkverbindung, Synchronisation und Verzögerungen führten dazu, dass die Steuerung von Pepper nicht in der gewünschten Präzision umgesetzt werden konnte. Diese technischen Herausforderungen erwiesen sich als besonders schwerwiegend und konnten trotz intensiver Bemühungen nicht vollständig überwunden werden.

\section{Komplexität der Integration}
Die Integration von Unity und ROS über den ROS-TCP-Connector stellte sich als komplexer als erwartet heraus. Obwohl Unity eine benutzerfreundliche Entwicklungsumgebung bietet, waren die Anforderungen an die ROS-Integration höher als erwartet, insbesondere im Hinblick auf Echtzeitfähigkeit und die nahtlose Zusammenarbeit der verschiedenen Softwarekomponenten. Die Komplexität dieser Integration erwies sich als Hindernis für die erfolgreiche Umsetzung des Projekts.

\section{Wahl von Unity}
Die Entscheidung für Unity als Entwicklungsplattform wurde aufgrund ihrer weiten Verbreitung, der umfangreichen Dokumentation und der großen Entwicklergemeinschaft getroffen. Unity bietet viele vorgefertigte Lösungen und Plugins, die die Entwicklung beschleunigen können. Dennoch stellte sich heraus, dass für diese spezielle Aufgabe tiefergehende Kenntnisse im Umgang mit der Unity-ROS-Integration erforderlich waren, als in der gegebenen Zeit erlernbar war. Diese Lernerfahrung zeigt deutlich die Bedeutung einer sorgfältigen Evaluierung und Auswahl der Entwicklungsplattform für komplexe Projekte.

\section{Potenzial für zukünftige Arbeiten}
Trotz der Herausforderungen zeigt das Projekt deutlich das Potenzial, das in der Kombination von VR und humanoider Robotik liegt. Künftige Arbeiten könnten von den gemachten Erfahrungen profitieren und die entwickelten Ansätze weiter verfolgen. Dies könnte die Erforschung alternativer Technologien für die VR-Steuerung, die Optimierung der Netzwerkverbindung und die Weiterentwicklung der Unity-ROS-Integration umfassen. Welche Schritte zur Weiterführung notwendig sind, werden auch im nachfolgenden Kapitel genauer erklärt.
\\

\noindent
Insgesamt war das Projekt eine wertvolle Lernerfahrung, die Einblicke in die komplexe Welt der Robotik und Virtual Reality ermöglichte. Die gesammelten Erfahrungen und Erkenntnisse werden als Grundlage für zukünftige Entwicklungen dienen, und wir sind zuversichtlich, dass die Vision einer nahtlosen Teleoperation von humanoiden Robotern über VR-Systeme mit weiterem Engagement und Forschung realisiert werden kann. Trotz des Fehlschlags des aktuellen Projekts sind wir motiviert, die gewonnenen Erkenntnisse in zukünftige Projekte zu integrieren und unser Wissen weiter auszubauen.

\chapter{Fortsetzung des Projekts}
Um das Projekt \textit{Pepper VR – Teleoperation eines humanoiden Roboters auf Basis der Analyse menschlicher Bewegung} in Zukunft erfolgreich umzusetzen, sind mehrere wichtige Schritte und Verbesserungen erforderlich. Nach der Analyse der bisherigen Herausforderungen und der gewonnenen Erkenntnisse haben wir folgende Empfehlungen für die Fortsetzung des Projekts zusammengestellt:

\section{Optimierung der technischen Infrastruktur}
Eine der größten Herausforderungen war die präzise und latenzfreie Datenübertragung zwischen der MetaQuest 3 VR-Brille und Pepper. Zur Verbesserung der technischen Infrastruktur sollten folgende Maßnahmen ergriffen werden:
\begin{itemize}
    \item \textbf{Verbesserung der Netzwerkverbindung}: Sicherstellen einer stabilen und schnellen Netzwerkverbindung, möglicherweise durch den Einsatz von dedizierten Netzwerken oder der Optimierung der bestehenden Infrastruktur. Dies kann durch die Verwendung von leistungsfähigeren Routern, Switches und Netzwerkkabeln erreicht werden.
    \item \textbf{Reduzierung der Latenzzeiten}: Implementierung von Echtzeit-Optimierungen sowohl auf der Hardware- als auch auf der Softwareseite, um die Verzögerungen bei der Datenübertragung zu minimieren. Dies kann durch die Optimierung von Netzwerkeinstellungen, die Auswahl von leistungsfähigeren Hardwarekomponenten und die Verwendung von effizienteren Datenübertragungsprotokollen erreicht werden.
\end{itemize}

\section{Erweiterung der Softwareintegration}
Die Integration von Unity und ROS über den ROS-TCP-Connector muss weiter verfeinert und genauer recherchiert werden. Hier sind die wesentlichen Schritte:
\begin{itemize}
    \item \textbf{Tiefere ROS-Integration}: Entwicklung zusätzlicher ROS-Nodes und -Services, um eine nahtlosere Kommunikation zwischen Unity und dem Robot Operating System zu gewährleisten. Dies kann durch die Entwicklung benutzerdefinierter ROS-Pakete und die Integration von ROS-Bibliotheken in Unity erreicht werden.
    \item \textbf{Fehlersuche und Debugging}: Intensive Fehlersuche und Debugging der bestehenden Implementierung, um die Ursachen für Synchronisationsprobleme zu identifizieren und zu beheben. Dies kann durch die Verwendung von Debugging-Tools wie ROS-Bag oder ROS-Log und die systematische Analyse von Protokollen und Nachrichtenverläufen erreicht werden.
\end{itemize}

\section{Testen und Validieren der VR-Steuerung}
Um die Steuerung von Pepper über die MetaQuest 3 VR-Brille erfolgreich umzusetzen, sind umfangreiche Tests und Validierungen notwendig:
\begin{itemize}
    \item \textbf{Erstellung von Testfällen}: Entwicklung spezifischer Testfälle, die die verschiedenen Aspekte der VR-Steuerung abdecken, um sicherzustellen, dass alle Funktionen wie erwartet arbeiten. Dies kann durch die Erstellung von automatisierten Testskripten und die systematische Durchführung von Benutzertests in verschiedenen Szenarien und Umgebungen erreicht werden.
    \item \textbf{Benutzerstudien}: Durchführung von Benutzerstudien, um die Benutzerfreundlichkeit und Effektivität der VR-Steuerung zu bewerten und basierend auf dem Feedback Verbesserungen vorzunehmen. Dies kann durch die Durchführung von Umfragen, Interviews und Beobachtungsstudien mit potenziellen Benutzern in realen Einsatzszenarien erreicht werden.
\end{itemize}

\section{Schulung und Dokumentation}
Die erfolgreiche Umsetzung des Projekts erfordert gut geschulte Teammitglieder und umfassende Dokumentation:
\begin{itemize}
    \item \textbf{Schulung der Teammitglieder}: Regelmäßige Schulungen und Workshops für das Team, um sicherzustellen, dass alle Mitglieder die notwendigen Kenntnisse und Fähigkeiten haben, um mit den verwendeten Technologien effektiv zu arbeiten. Dies kann durch die Organisation von Schulungsveranstaltungen, die Bereitstellung von Schulungsmaterialien und die Einrichtung von Online-Lernplattformen erreicht werden.
    \item \textbf{Erstellung umfassender Dokumentation}: Dokumentation aller Implementierungs- und Testprozesse, um eine klare und nachvollziehbare Basis für zukünftige Arbeiten zu schaffen. Dies kann durch die Erstellung von technischen Handbüchern, Benutzerhandbüchern, API-Dokumentationen und Projektberichten erreicht werden.
\end{itemize}

\section{Langfristige Wartung und Weiterentwicklung}
Um das Projekt langfristig erfolgreich zu halten, sollten regelmäßige Wartung und Weiterentwicklungen eingeplant werden:
\begin{itemize}
    \item \textbf{Regelmäßige Updates}: Kontinuierliche Aktualisierung der verwendeten Software und Hardware, um mit den neuesten Entwicklungen Schritt zu halten und die Stabilität und Sicherheit zu gewährleisten. Dies kann durch die regelmäßige Überprüfung und Aktualisierung von Bibliotheken, Frameworks und Betriebssystemen erreicht werden.
    \item \textbf{Weiterentwicklung der Funktionen}: Fortlaufende Erweiterung der Funktionen und Fähigkeiten des Systems, basierend auf den Bedürfnissen und Rückmeldungen der Benutzer. Dies kann durch die Implementierung neuer Features, die Integration von externen APIs und die Entwicklung von Erweiterungsmodulen erreicht werden.
\end{itemize}

\noindent
Durch die Umsetzung dieser Maßnahmen kann das Projekt \textit{Pepper VR} erfolgreich weitergeführt werden, um das volle Potenzial der Kombination von Virtual Reality und humanoider Robotik auszuschöpfen. Die gewonnenen Erkenntnisse und Erfahrungen bilden eine solide Grundlage für zukünftige Entwicklungen und Innovationen in diesem spannenden und zukunftsträchtigen Bereich.




% Ab hier beginnt der Anhang
\appendix
\addcontentsline{toc}{chapter}{Literaturverzeichnis}

% Haben Sie das "biblatex"-Paket nicht installiert, benutzen Sie folgendes:
% Ohne das "biblatex"-Paket (s. bericht.sty) produziert folgendes
% "deutsche" Zitate in Literaturverzeichnissen gemaß der Norm DIN 1505,
% Teil 2 vom Jan. 1984.
% Die Zitatmarken werden alphabetisch nach Verfassern
% sortiert und sind durch abgekürzte Verfasserbuchstaben plus
% Erscheinungsjahr in eckigen Klammern gekennzeichnet.

% \bibliographystyle{alphadin}
% \bibliography{bericht}

%%%%%%%%%%%%%%%%%%%%%%%%%%%%%%%%%%%%%%%5
% BIBLATEX
% Benutzt man das "biblatex"-Paket, muß man folgendes schreiben:
\def\refname{Literaturverzeichnis}
\printbibliography
%%%%%%%%%%%%%%%%%%%%%%%%%%%%%%%%%%%%%%%5


\end{document}
